\documentclass[11pt,a4paper]{article}

% ===================== POLICES & LANGUE =====================
\usepackage{fontspec}
\setmainfont{Latin Modern Roman}
\setsansfont{Latin Modern Sans}
\setmonofont{Latin Modern Mono}
\usepackage{polyglossia}
\setdefaultlanguage{french}

% ===================== MISE EN PAGE =====================
\usepackage[margin=2.5cm]{geometry}
\usepackage{amsmath,amssymb,graphicx,xcolor,booktabs,hyperref,tcolorbox,enumitem,fancyhdr,array,pifont,longtable,multirow,listings}
\tcbuselibrary{theorems,breakable,skins}
\usepackage{tikz}
\usetikzlibrary{shapes.geometric,arrows.meta,positioning,calc,backgrounds,fit}

\newcommand{\cmark}{\ding{51}}
\newcommand{\xmark}{\ding{55}}
\hypersetup{colorlinks=true,linkcolor=blue!60!black,urlcolor=blue!60!black,citecolor=blue!60!black}

% ===================== BOÎTES =====================
\newtcolorbox{objectifbox}[1][]{colback=blue!5,colframe=blue!50,fonttitle=\bfseries,title={#1},breakable}
\newtcolorbox{infobox}[1][]{colback=green!5,colframe=green!50!black,fonttitle=\bfseries,title={#1},breakable}
\newtcolorbox{attentionbox}[1][]{colback=orange!5,colframe=orange!70,fonttitle=\bfseries,title={#1},breakable}
\newtcolorbox{retenirbox}[1][]{colback=yellow!10,colframe=yellow!50!black,fonttitle=\bfseries,title={#1},breakable}
\newtcolorbox{prereqbox}[1][]{colback=cyan!5,colframe=cyan!50!black,fonttitle=\bfseries,title={#1},breakable}
\newtcolorbox{projetbox}[1][]{colback=teal!5,colframe=teal!50!black,fonttitle=\bfseries,title={#1},breakable}

% ===================== CODE (listings) =====================
\definecolor{codebg}{gray}{0.96}
\definecolor{kwcolor}{RGB}{0,90,160}
\definecolor{strcolor}{RGB}{160,55,30}
\definecolor{comcolor}{RGB}{95,120,95}
\definecolor{ndcolor}{RGB}{120,80,150}

\lstdefinelanguage{JS}{
  morekeywords={let,const,var,function,return,if,else,for,of,in,while,do,true,false,null,undefined,typeof,break,continue,new,this,async,await,try,catch,throw,finally},
  morekeywords=[2]{console,log,table,document,Math,window,alert,prompt,forEach,map,filter,push,pop,querySelector,querySelectorAll,addEventListener,classList,textContent,innerHTML,createElement,appendChild,getElementById,includes,fetch,json,JSON,then,stringify,parse,setItem,getItem,removeItem,setInterval,clearInterval},
  sensitive=true,
  morecomment=[l]{//},
  morecomment=[s]{/*}{*/},
  morestring=[b]",
  morestring=[b]',
  morestring=[b]`,
}

\lstdefinestyle{js}{
  language=JS,
  basicstyle=\ttfamily\small,
  keywordstyle=\color{kwcolor},
  keywordstyle=[2]\color{ndcolor},
  commentstyle=\color{comcolor}\itshape,
  stringstyle=\color{strcolor},
  backgroundcolor=\color{codebg},
  frame=single,
  rulecolor=\color{gray!45},
  breaklines=true,
  breakatwhitespace=false,
  showstringspaces=false,
  numbers=none,
  tabsize=2,
  columns=fullflexible,
  keepspaces=true,
  literate=%
    {é}{{é}}1 {è}{{è}}1 {à}{{à}}1 {â}{{â}}1 {ê}{{ê}}1
    {î}{{î}}1 {ô}{{ô}}1 {û}{{û}}1 {ù}{{ù}}1 {ç}{{ç}}1
    {ï}{{ï}}1 {ë}{{ë}}1 {É}{{É}}1 {È}{{È}}1 {À}{{À}}1
    {œ}{{œ}}1 {°}{{°}}1 {²}{{²}}1
}

\lstdefinestyle{html}{
  basicstyle=\ttfamily\small,
  backgroundcolor=\color{codebg},
  frame=single,
  rulecolor=\color{gray!45},
  breaklines=true,
  showstringspaces=false,
  numbers=none,
  tabsize=2,
  columns=fullflexible,
  keepspaces=true,
  literate=%
    {é}{{é}}1 {è}{{è}}1 {à}{{à}}1 {ê}{{ê}}1 {ç}{{ç}}1 {î}{{î}}1
}

\lstset{style=js}
\newcommand{\code}[1]{\texttt{#1}}

% ===================== EN-TÊTES =====================
\pagestyle{fancy}
\fancyhf{}
\fancyhead[L]{\small Jàng Afrig --- CM4 : Le Web moderne}
\fancyhead[R]{\small UCAD --- L1 DUT1-Info}
\fancyfoot[C]{\thepage}
\fancyfoot[L]{\small Dr. El Hadji Bassirou TOURÉ}
\fancyfoot[R]{\small 2025--2026}
\renewcommand{\headrulewidth}{0.4pt}
\renewcommand{\footrulewidth}{0.4pt}

% ===================== TITRE =====================
\title{%
\vspace{-1cm}
\textbf{\LARGE Fondamentaux du Web}\\
\textbf{\LARGE HTML, CSS, JavaScript}\\[0.8cm]
\textbf{\huge CM4 --- Le Web moderne}\\[0.3cm]
\textbf{\Large APIs, JSON et code asynchrone}\\[0.5cm]
\large Licence 1 DUT1-Informatique --- 2025--2026}
\date{}
\author{}

\begin{document}
\maketitle
\thispagestyle{empty}

\begin{prereqbox}[Avant ce cours]
Vous avez construit, au Lab 7, un quiz jouable\ldots\ mais limité à dix pays écrits à la main. Ce cours explique comment une page web va chercher de \emph{vraies} données en ligne, automatiquement. C'est ce qui permettra à Jàng Afrig de connaître les 54 pays africains sans que vous ayez à les taper.

\smallskip
\textbf{Prérequis :} les fonctions et les objets (CM3, Labs 5--6), et le quiz du Lab 7.
\end{prereqbox}

\begin{objectifbox}[Objectifs du CM4]
À la fin de ce cours, vous comprendrez :
\begin{enumerate}[leftmargin=*,noitemsep]
  \item Ce qu'est une \textbf{API} et le style \textbf{REST}.
  \item Le format \textbf{JSON}, langage d'échange du web.
  \item La différence entre code \textbf{synchrone} et \textbf{asynchrone}.
  \item Comment faire une requête avec \code{fetch} + \code{async}/\code{await}.
  \item Comment se protéger des erreurs avec \code{try}/\code{catch}.
  \item L'API \textbf{REST Countries}, que nous utiliserons dès le Lab 8.
\end{enumerate}
\end{objectifbox}

\tableofcontents
\bigskip\hrule\bigskip

% ===================== SECTION 1 =====================
\section{Qu'est-ce qu'une API ?}

Jusqu'ici, Jàng Afrig vivait en autarcie : toutes ses données étaient dans \code{data.js}. Mais la plupart des applications ont besoin de données qui vivent \emph{ailleurs}, sur un serveur distant : la météo, les prix, la liste des pays\ldots\ Comment les obtenir ?

Grâce à une \textbf{API} (\emph{Application Programming Interface}). Une API, c'est une \textbf{porte d'entrée} qu'un service met à disposition pour qu'on lui demande des données de façon organisée.

\begin{infobox}[L'analogie du restaurant]
Imaginez un restaurant :
\begin{itemize}[leftmargin=*,noitemsep]
  \item \textbf{Vous}, à table, êtes la \emph{page web}.
  \item La \textbf{cuisine} (que vous ne voyez pas) est la \emph{base de données} du service.
  \item Le \textbf{serveur} qui fait la navette est l'\emph{API}.
  \item Le \textbf{menu} liste ce que vous pouvez commander (les \emph{adresses} disponibles).
  \item Votre \textbf{commande} est la \emph{requête} ; le \textbf{plat} qui arrive est la \emph{réponse}.
\end{itemize}
Vous n'entrez jamais dans la cuisine : vous passez \emph{toujours} par le serveur. De même, une page ne touche jamais directement la base de données : elle passe par l'API.
\end{infobox}

% ===================== SECTION 2 =====================
\section{REST : commander proprement}

\textbf{REST} est le style d'API le plus répandu sur le web. Son idée : chaque donnée a une \textbf{adresse} (une URL), et on agit dessus avec un petit nombre de \textbf{verbes}.

\begin{center}
\begin{tabular}{l l l}
\toprule
\textbf{Verbe HTTP} & \textbf{Action} & \textbf{Au restaurant} \\
\midrule
\code{GET}    & Lire une donnée      & Consulter le menu / recevoir un plat \\
\code{POST}   & Créer une donnée     & Passer une nouvelle commande \\
\code{PUT}    & Modifier une donnée  & Changer sa commande \\
\code{DELETE} & Supprimer une donnée & Annuler sa commande \\
\bottomrule
\end{tabular}
\end{center}

\begin{retenirbox}[Pour Jàng Afrig, un seul verbe]
Notre quiz ne fait que \textbf{lire} des données (la liste des pays). Nous n'utiliserons donc que le verbe \code{GET} --- le plus simple. C'est d'ailleurs ce que fait votre navigateur quand vous tapez une adresse : un \code{GET}.
\end{retenirbox}

% ===================== SECTION 3 =====================
\section{JSON : la langue d'échange}

Quand l'API répond, elle n'envoie pas une page web : elle envoie des \textbf{données brutes}, dans un format que tout le monde comprend, le \textbf{JSON} (\emph{JavaScript Object Notation}).

Bonne nouvelle : le JSON ressemble \emph{trait pour trait} aux objets et tableaux JavaScript que vous connaissez déjà (CM3, Lab 6).

\begin{lstlisting}[style=js]
{
  "cca3": "SEN",
  "name": { "common": "Senegal" },
  "capital": ["Dakar"],
  "languages": { "fra": "French", "wol": "Wolof" }
}
\end{lstlisting}

On y reconnaît : des \textbf{objets} (entre \code{\{ \}}), des \textbf{tableaux} (entre \code{[ ]}), du texte, des nombres. La seule particularité : en JSON, les noms de propriétés sont \emph{toujours} entre guillemets.

\begin{attentionbox}[Attention à la forme des champs]
Deux pièges fréquents avec REST Countries :
\begin{itemize}[leftmargin=*,noitemsep]
  \item Le nom est \emph{niché} : on y accède par \code{pays.name.common} (et non \code{pays.name}).
  \item La capitale est un \textbf{tableau} : on prend le premier élément, \code{pays.capital[0]}.
\end{itemize}
\end{attentionbox}

% ===================== SECTION 4 =====================
\section{Synchrone ou asynchrone : pourquoi Internet est « lent »}

Aller chercher des données en ligne prend du temps (le signal voyage, le serveur réfléchit). Pendant ce temps, que fait votre page ? Deux stratégies :

\begin{infobox}[L'analogie du café]
\textbf{Au comptoir (synchrone) :} vous commandez un café, et vous \emph{restez planté} devant la machine jusqu'à ce qu'il soit prêt. Vous ne faites rien d'autre. Si la page faisait ça, elle se \textbf{figerait} à chaque requête.

\smallskip
\textbf{Via une application (asynchrone) :} vous commandez, puis vous allez vous asseoir et discuter. Quand le café est prêt, une \textbf{notification} vous prévient. La page reste \emph{vivante} pendant l'attente.
\end{infobox}

Le web moderne est \textbf{asynchrone} : on lance la requête, la page continue de fonctionner, et on traite la réponse \emph{quand elle arrive}.

\subsection{Les promesses}
Quand on lance une requête asynchrone, on ne reçoit pas la donnée tout de suite : on reçoit une \textbf{promesse} (\emph{promise}), c'est-à-dire un objet qui dit \emph{« je te répondrai plus tard, avec le résultat ou avec une erreur »}. C'est le ticket que vous donne le café : la preuve qu'une réponse arrivera.

% ===================== SECTION 5 =====================
\section{\texttt{fetch} + \texttt{async}/\texttt{await} : la requête en pratique}

\code{fetch(url)} lance une requête vers une adresse. Comme c'est asynchrone, on utilise deux mots-clés pour attendre proprement la réponse :
\begin{itemize}[leftmargin=*,noitemsep]
  \item \code{async} devant une fonction : « cette fonction contient de l'attente ».
  \item \code{await} devant une opération : « attends ici que ce soit prêt, sans figer la page ».
\end{itemize}

\begin{lstlisting}[style=js]
async function chargerPays() {
  const reponse = await fetch("https://restcountries.com/v3.1/region/africa");
  const donnees = await reponse.json(); // transformer la réponse en objets JS
  return donnees;
}
\end{lstlisting}

\textbf{Pourquoi deux \code{await} ?}
\begin{enumerate}[leftmargin=*,noitemsep]
  \item Le premier attend que le serveur \emph{réponde} (le facteur revient avec le colis).
  \item Le second attend que la réponse soit \emph{convertie} en objets JavaScript utilisables, avec \code{.json()} (on ouvre le colis).
\end{enumerate}

\begin{center}
\begin{tikzpicture}[
  font=\footnotesize,
  boite/.style={rectangle, rounded corners=4pt, draw, very thick, align=center, minimum height=1.3cm, minimum width=3.2cm},
  arr/.style={-{Stealth[length=3mm]}, very thick},
  lab/.style={font=\scriptsize, align=center}
]
  \node[boite, draw=blue!55, fill=blue!6] (client) {\textbf{Votre page}\\\scriptsize (le navigateur)};
  \node[boite, draw=teal!55!black, fill=teal!6, right=5.2cm of client] (api) {\textbf{API REST Countries}\\\scriptsize (serveur distant)};

  \draw[arr, blue!55!black] ([yshift=4mm]client.east) -- node[lab, above]{\textbf{1.} requête \code{GET}\\\scriptsize \code{fetch(url)}} ([yshift=4mm]api.west);
  \draw[arr, teal!55!black] ([yshift=-4mm]api.west) -- node[lab, below]{\textbf{2.} réponse\\\scriptsize données \code{JSON}} ([yshift=-4mm]client.east);
\end{tikzpicture}

\textit{Le dialogue client $\leftrightarrow$ API : la page demande, l'API répond en JSON.}
\end{center}

% ===================== SECTION 6 =====================
\section{\texttt{try}/\texttt{catch} : le filet de sécurité}

Et si la connexion saute, ou si le serveur ne répond pas ? Sans précaution, l'application plante. On entoure donc la requête d'un \textbf{filet} : \code{try} (essaie) / \code{catch} (attrape l'erreur).

\begin{lstlisting}[style=js]
async function chargerPays() {
  try {
    const reponse = await fetch("https://restcountries.com/v3.1/region/africa");
    const donnees = await reponse.json();
    return donnees;
  } catch (erreur) {
    console.log("Connexion impossible :", erreur);
    return []; // on renvoie une liste vide plutôt que de planter
  }
}
\end{lstlisting}

\emph{Analogie :} c'est la ceinture de sécurité. On espère ne jamais s'en servir, mais le jour où ça secoue, elle évite la catastrophe.

% ===================== SECTION 7 =====================
\section{Notre fournisseur : REST Countries v3.1}

C'est l'API que nous utiliserons : \textbf{gratuite, sans inscription, sans clé}. Pour récupérer les pays d'Afrique, on demande la région \code{africa}, en précisant les \textbf{champs} qui nous intéressent (pour ne recevoir que l'utile) :

\begin{lstlisting}[style=html]
https://restcountries.com/v3.1/region/africa?fields=name,capital,cca3,flags,languages,borders
\end{lstlisting}

\begin{center}
\begin{tabular}{l l}
\toprule
\textbf{Champ} & \textbf{Contenu (et accès)} \\
\midrule
\code{name.common}  & Nom courant du pays --- \code{pays.name.common} \\
\code{capital}      & Capitale (tableau) --- \code{pays.capital[0]} \\
\code{cca3}         & Code à 3 lettres --- \code{pays.cca3} (ex. \code{"SEN"}) \\
\code{flags.svg}    & Adresse de l'image du drapeau \\
\code{languages}    & Langues parlées (objet) \\
\code{borders}      & Pays voisins (tableau de codes) \\
\bottomrule
\end{tabular}
\end{center}

\begin{infobox}[Démo en direct]
Collez l'adresse ci-dessus dans la barre d'un navigateur : vous verrez le JSON brut d'une cinquantaine de pays africains. Puis, dans la console (F12), nous lancerons \code{chargerPays()} et nous afficherons le résultat avec \code{console.table} : les vrais pays apparaissent, en direct.
\end{infobox}

\begin{retenirbox}[Les réflexes du jour]
\begin{enumerate}[leftmargin=*,noitemsep]
  \item Une API se demande avec une \textbf{URL} et un \textbf{verbe} (pour nous : \code{GET}).
  \item La réponse arrive en \textbf{JSON} (objets et tableaux familiers).
  \item On attend la réponse avec \code{await fetch(...)} puis \code{await reponse.json()}.
  \item On protège toujours la requête avec \code{try}/\code{catch}.
\end{enumerate}
\end{retenirbox}

% ===================== MINI-EXO =====================
\section{Mini-exercice en classe}

\begin{objectifbox}[À faire ensemble, dans la console]
\begin{enumerate}[leftmargin=*]
  \item Ouvrez une page quelconque, puis la console (F12).
  \item Copiez-collez la fonction \code{chargerPays} (version avec \code{try}/\code{catch}).
  \item Tapez : \code{const pays = await chargerPays();} puis \code{console.table(pays);}.
  \item Combien de pays voyez-vous ? Affichez le nom courant du premier : \code{pays[0].name.common}.
\end{enumerate}
\emph{Astuce :} la console moderne accepte \code{await} directement, sans fonction autour.
\end{objectifbox}

% ===================== RÉSUMÉ =====================
\section{Résumé et transition}

Une API est une porte d'entrée vers des données distantes ; on la sollicite avec une URL et le verbe \code{GET} ; elle répond en JSON ; et on récupère ce JSON sans figer la page grâce à \code{async}/\code{await}, le tout protégé par \code{try}/\code{catch}.

Au \textbf{Lab 8}, on met tout cela en pratique : on remplace les dix pays écrits à la main par les \textbf{vrais pays africains} de REST Countries. Le mode Capitales continuera de fonctionner --- mais avec, cette fois, toute l'Afrique.

\begin{center}
\textit{La théorie est posée. Au Lab 8, Jàng Afrig se connecte au monde réel.}
\end{center}

\bigskip
\noindent\rule{\textwidth}{0.4pt}
\begin{center}
\small Dr.~El Hadji Bassirou TOURÉ --- UCAD\\
\small Licence 1 DUT1-Informatique --- Fondamentaux du Web --- 2025--2026
\end{center}

\end{document}
